\newpage
\section{Comparaci\'on de centralidades}

\subsection{Minired del S\&P 500}


Notemos que en HITS, el algoritmo predeterminado se quedo con la componente que
ten\'ia mayor peso espectral. Eso no es un error, pero es importante aclarar que
las puntuaciones de \textit{hubs} y \textit{authorities} se calculan solo para esa
parte del grafo. En cambio, el \textit{degree} se calcula para toda la red,
incluyendo nodos aislados o componentes menores. Por eso es posible que algunos
nodos tengan \textit{degree} positivo pero puntuaciones de HITS iguales a cero,
porque quedaron fuera de la componente principal que se us\'o para el c\'alculo
espectral.

\sidefigure{figs/spideg.png}{Distribuci\'on del \textit{in-degree} en la minired del S\&P 500.}{fig:spideg}{%
Este histograma muestra qu\'e tan concentrado est\'a el lado comprador. La mayor\'ia de las empresas recibe de muy pocos proveedores, Destacan \text{WMT} y \text{PLTR}, con cuatro proveedores cada una, y justo detr\'as aparecen \text{AAPL} y \text{AMZN}, con tres. La interpretaci\'on econ\'omica cambia seg\'un el nodo: \text{WMT} junta consumo y pagos; \text{PLTR} recibe de plataformas tecnol\'ogicas; \text{AMZN} combina compras tecnol\'ogicas con su propio papel de distribuci\'on hacia otros sectores.%
}

\sidefigure{figs/spodeg.png}{Distribuci\'on del \textit{out-degree} en la minired del S\&P 500.}{fig:spodeg}{%
En \textit{out-degree} \text{NVDA} sobresale con mucha claridad, porque vende a siete nodos del grafo; luego vienen \text{AMZN}, \text{AMD}, \text{MSFT} y \text{MA}. Hay pocos proveedores que vendan a muchos compradores y hay una mayor\'ia de firmas con alcance corto. Esta figura tambi\'en aclara por qu\'e \text{MU} no sale como gran vendedor en extensi\'on: en esta minired solo tiene una salida, hacia \text{NVDA}. 
}


\sidefigure{figs/sphubs.png}{Comparaci\'on de centralidad de \textit{hubs} en la minired del S\&P 500.}{fig:sphubs}{%
En esta gr\'afica los \textit{hubs} fuertes son \text{NVDA} y \text{AMD}, la arista se entiende como proveedor $\to$ comprador y ambas aparecen como firmas que abastecen a varias empresas importantes. \text{NVDA} domina porque le vende a \text{MSFT}, \text{AMZN}, \text{GOOG}, \text{GOOGL}, \text{ORCL}, \text{META} y \text{TSLA}; es decir, no solo vende mucho, sino que vende a nodos que ya ocupan una posici\'on fuerte en la parte compradora de la red. \text{MU}, en cambio, queda mucho m\'as arriba de la cadena que del centro del grafo: aqu\'i solo le vende a \text{NVDA}, as\'i que su papel se ve m\'as como proveedor especializado que como vendedor amplio.%
}

\sidefigure{figs/spauth.png}{Comparaci\'on de centralidad de \textit{authorities} en la minired del S\&P 500.}{fig:spauth}{%
Las \textit{authorities} cuentan la historia del lado comprador. Aqu\'i suben \text{AMZN}, \text{MSFT}, \text{TSLA} y \text{ORCL} porque reciben de proveedores que ya eran fuertes como \textit{hubs}, sobre todo \text{NVDA} y \text{AMD}. \text{WMT} s\'i es una de las firmas que m\'as compra si se miran solo enlaces entrantes, porque recibe de \text{AAPL}, \text{PG}, \text{BRK.B} y \text{MA}. Pero HITS no premia solamente cu\'antos proveedores tienes, sino de qui\'en vienen esos enlaces y en qu\'e parte del grafo est\'as parado. Por eso \text{WMT} aparece como comprador local fuerte, aunque no como autoridad dominante en sentido espectral.%
}


La comparaci\'on hace evidente que hay 3 roles distintos. El primero es el del proveedor y ah\'i \text{NVDA} es el caso m\'as claro. El segundo es el del comprador, como \text{AMZN} y \text{MSFT}, porque reciben de nodos muy influyentes. El tercero es el del intermediario. Ah\'i la betweenness resulta especialmente \'util: \text{AMZN} no sobresale como \textit{hub}, pero s\'i como puente entre el bloque de proveedores y clientes como \text{NFLX}, \text{XOM}, \text{JNJ}, \text{ABBV} o \text{PLTR}. \text{MSFT} cumple algo parecido hacia \text{CVX}, \text{GE}, \text{UNH} y tambi\'en \text{PLTR}. 

Tambi\'en conviene separar \textit{degree} de HITS. Por \textit{out-degree}, \text{AMZN} vende a cinco nodos y \text{MA} a cuatro, pero eso no basta para convertirlas en \textit{hubs} dominantes, porque muchas de esas ventas terminan en nodos hoja o en componentes con menos peso espectral. En HITS importa que tan valiosos son los destinos de esos enlaces. Del mismo modo, por \textit{in-degree} \text{PLTR} y \text{WMT} son compradores grandes, pero la autoridad fuerte se concentra en \text{NVDA} y \text{AMD}. 

Tomados juntos, los dos histogramas muestran una red asim\'etrica, no son los mismos en entradas y salidas. En la minired la importancia depende de si una empresa abastece a muchas otras, si compra de muchas o si hace ambas cosas en zonas distintas del grafo. 

\subsection{Red de \textit{Star Wars Episode II}}

En una red de coapariciones se puede observar que los personajes principales son \text{PADM\'E}, \text{OBI-WAN} y \text{ANAKIN}. \text{PADM\'E} concentra la subtrama pol\'itica y la relaci\'on con Anakin; \text{OBI-WAN} la subtrama sobre la investigaci\'on; y \text{ANAKIN} es el principal.


\sidefigure{figs/sweigv.png}{Comparaci\'on de centralidad espectral en la red de coapariciones de \textit{Star Wars Episode II}.}{fig:sweigv}{%
La centralidad espectral sube a quienes no solo aparecen mucho, sino que adem\'as est\'an conectados con otros personajes tambi\'en centrales. Por eso \text{PADM\'E} suele quedar ligeramente arriba: pasa por el Senado, por Naboo, por su familia y por la trama con \text{ANAKIN}. \text{OBI-WAN} aparece enseguida porque enlaza el Consejo Jedi, a \text{Dexter}, a \text{Kamino}, a \text{Jango Fett} y a \text{Geonosis}. \text{ANAKIN} queda muy cerca, pero buena parte de su recorrido est\'a pegado a \text{PADM\'E}; por eso su papel es muy central sin abrir tantas ramas nuevas como el de Obi-Wan.%
}

\sidefigure{figs/swbtw.png}{Centralidad de intermediaci\'on en la red de coapariciones de \textit{Star Wars Episode II}.}{fig:swbtw}{%
Aqu\'i se ve mejor qu\'e \text{OBI-WAN} termina arriba, empieza con el atentado y el Consejo Jedi, sigue el rastro hasta \text{Kamino}, persigue a \text{Jango Fett}, llega a \text{Geonosis} y acaba frente a \text{Count Dooku}. \text{PADM\'E} tambi\'en intermedia bastante porque une el bloque pol\'itico de Coruscant, el regreso a Naboo y la arena final. \text{ANAKIN} es central, pero pasa tramos m\'as largos en escenas de pareja o de familia; por eso intermedia menos que ambos.%
}


\begin{figure}[H]
    \centering
    \begin{subfigure}[t]{0.48\textwidth}
        \centering
        \reportincludegraphics{\linewidth}{figs/swatuh.png}
        \caption{Comparaci\'on entre \textit{authorities} y \textit{hubs} en la red de \textit{Star Wars Episode II}.}
        \label{fig:swatuh}
    \end{subfigure}\hfill
    \begin{subfigure}[t]{0.48\textwidth}
        \centering
        \reportincludegraphics{\linewidth}{figs/swhubs.png}
        \caption{Comparaci\'on de centralidades tipo HITS en la red de coapariciones de \textit{Star Wars Episode II}.}
        \label{fig:swhubs}
    \end{subfigure}
\end{figure}


No hay diferencia entre \textit{authorities} y \textit{hubs}, y eso es exactamente lo que uno esperar\'ia en una red no dirigida. No tiene sentido separar al que ``apunta'' del que ``es apuntado'', porque nadie le vende ni le compra nada a otro personaje



 \textit{Episode II} se trata tres temas: la intriga pol\'itica alrededor de \text{Padm\'e}, la investigaci\'on Jedi de \text{Obi-Wan} y la historia de \text{Anakin} con Padm\'e y con su pasado en Tatooine. Por eso los tres dominan casi todas las m\'etricas. Detr\'as aparecen \text{Yoda} y \text{Mace Windu}, que suben por el peso del Consejo Jedi y la batalla final, y luego personajes como \text{Jar Jar}, \text{Palpatine} o \text{Mas Amedda}, que se benefician de estar cerca del bloque pol\'itico. 